%═══════════════════════════════════════════════════════════════════════════════
\chapter{Fundamentos de Probabilidad y Estadística}
\label{ap:probabilidad}
%═══════════════════════════════════════════════════════════════════════════════

Este apéndice recoge los conceptos de probabilidad que el libro usa pero que no
forman parte del currículo estándar de física de pregrado. No es un repaso de
probabilidad básica: la gaussiana, el teorema de Bayes y el teorema central del
límite se dan por conocidos. El objetivo es preciso: proporcionar el vocabulario
y los resultados que hacen legibles los Capítulos~\ref{ch:cap7},
\ref{ch:cap10} y~\ref{ch:cap11}.

Cada sección indica explícitamente dónde se usa en el libro.

%───────────────────────────────────────────────────────────────────────────────
\section{Notación y repaso de referencia}
\label{sec:ap_notacion}
%───────────────────────────────────────────────────────────────────────────────

Se usa la siguiente notación a lo largo del libro. Sea $(\Omega, \mathcal{F},
\mathbb{P})$ un espacio de probabilidad, donde $\Omega$ es el espacio muestral,
$\mathcal{F}$ una $\sigma$-álgebra de eventos, y $\mathbb{P}$ una medida de
probabilidad. Una \textbf{variable aleatoria} $X: \Omega \to \mathbb{R}$ es una
función medible. Su \textbf{densidad} (cuando existe) $p_X(x)$ satisface
$\mathbb{P}(X \in A) = \int_A p_X(x)\,dx$.

\medskip
La siguiente tabla recoge las distribuciones que aparecen con frecuencia en el
texto. Los momentos se dan sin derivación.

\begin{center}
\begin{tabular}{lllll}
\toprule
\textbf{Distribución} & \textbf{Parámetros} & \textbf{Densidad $p(x)$}
  & $\mathbb{E}[X]$ & $\mathrm{Var}[X]$ \\
\midrule
Gaussiana & $\mu\in\mathbb{R}$, $\sigma^2>0$
  & $\frac{1}{\sqrt{2\pi\sigma^2}}e^{-(x-\mu)^2/(2\sigma^2)}$
  & $\mu$ & $\sigma^2$ \\[4pt]
Gaussiana mult. & $\bm{\mu}\in\mathbb{R}^d$, $\bm{\Sigma}\succ 0$
  & $(2\pi)^{-d/2}|\bm{\Sigma}|^{-1/2}e^{-\frac{1}{2}(\bm{x}-\bm{\mu})^\top\bm{\Sigma}^{-1}(\bm{x}-\bm{\mu})}$
  & $\bm{\mu}$ & $\bm{\Sigma}$ \\[4pt]
Gibbs-Boltzmann & $U:\mathbb{R}^d\to\mathbb{R}$, $T>0$
  & $Z^{-1}e^{-U(\bm{x})/(k_BT)}$
  & --- & --- \\[4pt]
Bernoulli & $p\in[0,1]$
  & $p^x(1-p)^{1-x}$, $x\in\{0,1\}$
  & $p$ & $p(1-p)$ \\[4pt]
Poisson & $\lambda>0$
  & $e^{-\lambda}\lambda^x/x!$, $x\in\mathbb{N}_0$
  & $\lambda$ & $\lambda$ \\
\bottomrule
\end{tabular}
\end{center}

\medskip
La \textbf{gaussiana multivariada} merece atención especial porque es el objeto
central de los Capítulos~\ref{ch:cap10}--\ref{ch:cap12}. Sus propiedades clave:
\begin{itemize}
  \item Si $\bm{X}\sim\mathcal{N}(\bm{\mu},\bm{\Sigma})$ y $\bm{A}$ es una
        matriz constante, entonces $\bm{A}\bm{X}\sim\mathcal{N}(\bm{A}\bm{\mu},
        \bm{A}\bm{\Sigma}\bm{A}^\top)$.
  \item Si $\bm{X}\sim\mathcal{N}(\bm{\mu},\bm{\Sigma})$, la distribución
        condicional $\bm{X}_1|\bm{X}_2 = \bm{x}_2$ es también gaussiana, con
        media $\bm{\mu}_1 + \bm{\Sigma}_{12}\bm{\Sigma}_{22}^{-1}(\bm{x}_2 -
        \bm{\mu}_2)$ y covarianza $\bm{\Sigma}_{11} - \bm{\Sigma}_{12}
        \bm{\Sigma}_{22}^{-1}\bm{\Sigma}_{21}$ (fórmula de Schur).
  \item La suma de gaussianas independientes es gaussiana:
        $X\sim\mathcal{N}(\mu_1,\sigma_1^2)$, $Y\sim\mathcal{N}(\mu_2,\sigma_2^2)$
        independientes $\Rightarrow$ $X+Y\sim\mathcal{N}(\mu_1+\mu_2,
        \sigma_1^2+\sigma_2^2)$.
\end{itemize}

%───────────────────────────────────────────────────────────────────────────────
\section{Probabilidad bayesiana y divergencia KL}
\label{sec:ap_bayesiana}
%───────────────────────────────────────────────────────────────────────────────

\textit{Usado en: Capítulo~\ref{ch:cap7} (inferencia bayesiana, ELBO),
Capítulo~\ref{ch:cap11} (score matching), Capítulo~\ref{ch:cap12} (CFG).}

\medskip
La regla de Bayes en notación de densidades:
\[
  p(\bm{\theta}|\bm{x}) = \frac{p(\bm{x}|\bm{\theta})\,p(\bm{\theta})}{p(\bm{x})},
\]
donde $p(\bm{\theta})$ es el \textbf{prior}, $p(\bm{x}|\bm{\theta})$ la
\textbf{verosimilitud}, $p(\bm{\theta}|\bm{x})$ el \textbf{posterior}, y
$p(\bm{x}) = \int p(\bm{x}|\bm{\theta})p(\bm{\theta})\,d\bm{\theta}$ la
\textbf{evidencia} (constante de normalización respecto a $\bm{\theta}$).

\subsection{La divergencia de Kullback-Leibler}
\label{subsec:ap_kl}

\begin{definicion}{Divergencia KL}
\label{def:ap_kl}
Para dos densidades $p$ y $q$ sobre $\mathbb{R}^d$:
\begin{equation}
  D_{\mathrm{KL}}(p\|q)
  = \int p(\bm{x})\log\frac{p(\bm{x})}{q(\bm{x})}\,d\bm{x}
  = \mathbb{E}_p\!\left[\log\frac{p(\bm{x})}{q(\bm{x})}\right].
  \label{eq:ap_kl}
\end{equation}
\end{definicion}

La divergencia KL no es una distancia: en general $D_{\mathrm{KL}}(p\|q) \neq
D_{\mathrm{KL}}(q\|p)$. Sus propiedades fundamentales:

\begin{proposicion}{Propiedades de la KL}
\label{prop:ap_kl}
\begin{enumerate}
  \item \textbf{No negatividad:} $D_{\mathrm{KL}}(p\|q) \geq 0$, con igualdad
        si y solo si $p = q$ casi en todas partes.
  \item \textbf{Desigualdad de Jensen:} para cualquier función convexa $\phi$
        y variable aleatoria $X$: $\phi(\mathbb{E}[X]) \leq \mathbb{E}[\phi(X)]$.
        La no negatividad de la KL es un caso especial con $\phi = -\log$.
  \item \textbf{Aditiva bajo independencia:} si $p(\bm{x},\bm{y}) =
        p_1(\bm{x})p_2(\bm{y})$ y $q(\bm{x},\bm{y}) = q_1(\bm{x})q_2(\bm{y})$,
        entonces $D_{\mathrm{KL}}(p\|q) = D_{\mathrm{KL}}(p_1\|q_1) +
        D_{\mathrm{KL}}(p_2\|q_2)$.
\end{enumerate}
\end{proposicion}

\begin{proof}
(1) Por la desigualdad de Jensen con $\phi = -\log$ (convexa):
$D_{\mathrm{KL}}(p\|q) = -\mathbb{E}_p[\log(q/p)] \geq -\log\mathbb{E}_p[q/p]
= -\log\int q\,d\bm{x} = 0$.
La igualdad requiere $q/p$ constante c.s., que junto con $\int p = \int q = 1$
da $p = q$. (3) Directo desde la definición y la aditividad del logaritmo.
$\square$
\end{proof}

\subsection{La ELBO como consecuencia de la KL}
\label{subsec:ap_elbo}

La \textbf{cota inferior variacional} (ELBO) es la herramienta central del
Capítulo~\ref{ch:cap7} y de DDPM (Capítulo~\ref{ch:cap12}). Se obtiene
directamente de la no negatividad de la KL.

Sea $p_\theta(\bm{x}) = \int p_\theta(\bm{x},\bm{z})\,d\bm{z}$ una distribución
con variables latentes $\bm{z}$ intractable. Para cualquier distribución
variacional $q_\phi(\bm{z}|\bm{x})$:
\begin{align}
  \log p_\theta(\bm{x})
  &= \log\int p_\theta(\bm{x},\bm{z})\,d\bm{z} \notag \\
  &= \log\int q_\phi(\bm{z}|\bm{x})\frac{p_\theta(\bm{x},\bm{z})}
    {q_\phi(\bm{z}|\bm{x})}\,d\bm{z} \notag \\
  &\geq \mathbb{E}_{q_\phi}\!\left[\log\frac{p_\theta(\bm{x},\bm{z})}
    {q_\phi(\bm{z}|\bm{x})}\right]
  =: \mathrm{ELBO}(\theta,\phi).
  \label{eq:ap_elbo}
\end{align}
El paso de la segunda a la tercera línea es la desigualdad de Jensen. La brecha
es exactamente la KL:
\begin{equation}
  \log p_\theta(\bm{x}) = \mathrm{ELBO}(\theta,\phi)
  + D_{\mathrm{KL}}(q_\phi(\bm{z}|\bm{x})\|p_\theta(\bm{z}|\bm{x})).
  \label{eq:ap_elbo_gap}
\end{equation}
Maximizar el ELBO respecto a $\phi$ minimiza $D_{\mathrm{KL}}(q_\phi\|p_\theta)$,
acercando la distribución variacional al posterior verdadero.

\subsection{Información de Fisher}
\label{subsec:ap_fisher}

\textit{Usado en: Capítulo~\ref{ch:cap7} (gradiente natural, cota de Cramér-Rao),
Capítulo~\ref{ch:cap11} (H-teorema, tasa de decrecimiento de la KL).}

\begin{definicion}{Información de Fisher}
\label{def:ap_fisher}
Para una familia paramétrica $\{p_\theta\}$, la \textbf{matriz de información
de Fisher} es:
\begin{equation}
  \bm{I}(\theta) = \mathbb{E}_{p_\theta}\!\left[
    \nabla_\theta\log p_\theta(\bm{x})\,
    \nabla_\theta\log p_\theta(\bm{x})^\top
  \right]
  = -\mathbb{E}_{p_\theta}\!\left[\nabla^2_\theta\log p_\theta(\bm{x})\right].
  \label{eq:ap_fisher}
\end{equation}
\end{definicion}

La igualdad entre las dos expresiones de~\eqref{eq:ap_fisher} se verifica
diferenciando la identidad $\int p_\theta\,d\bm{x} = 1$ dos veces respecto
a $\theta$.

La información de Fisher cuantifica cuánto varía la distribución al cambiar
$\theta$: es la curvatura de la log-verosimilitud. Su interpretación estadística
es la \textbf{cota de Cramér-Rao}: para cualquier estimador insesgado
$\hat{\theta}$ de $\theta$, $\mathrm{Var}[\hat{\theta}] \geq \bm{I}(\theta)^{-1}$.

En el Capítulo~\ref{ch:cap11} aparece en un contexto diferente: como la tasa de
disipación de la divergencia KL hacia el equilibrio. Para un proceso con balance
detallado respecto a $p^*$:
\[
  \frac{d}{dt}D_{\mathrm{KL}}(p_t\|p^*)
  = -\frac{\sigma^2}{2}\,\mathcal{I}(p_t\|p^*),
\]
donde $\mathcal{I}(p_t\|p^*) = \mathbb{E}_{p_t}[\|\nabla\log p_t -
\nabla\log p^*\|^2]$ es la \textbf{información de Fisher relativa}. Esta
identidad es el H-teorema de Boltzmann en forma cuantitativa.

%───────────────────────────────────────────────────────────────────────────────
\section{Procesos estocásticos}
\label{sec:ap_procesos}
%───────────────────────────────────────────────────────────────────────────────

\textit{Usado en: Capítulo~\ref{ch:cap10} (movimiento browniano, integral de
Itô), Capítulo~\ref{ch:cap11} (Fokker-Planck, proceso reverso).}

\medskip
Un \textbf{proceso estocástico} es una familia de variables aleatorias
$\{X_t\}_{t\in T}$ indexada por el tiempo $T\subseteq[0,\infty)$, todas
definidas sobre el mismo espacio de probabilidad $(\Omega,\mathcal{F},\mathbb{P})$.
Para cada $\omega\in\Omega$ fijo, la función $t\mapsto X_t(\omega)$ es una
\textbf{trayectoria} o \textbf{realización} del proceso.

En el libro se trabaja con tiempo continuo $T=[0,T_{\max}]$. La distribución
del proceso está completamente determinada por sus
\textbf{distribuciones finito-dimensionales}: para cualquier
$t_1 < t_2 < \cdots < t_n$, la ley conjunta de $(X_{t_1},\ldots,X_{t_n})$.

\subsection{Filtración y procesos adaptados}
\label{subsec:ap_filtracion}

La \textbf{filtración} $\{\mathcal{F}_t\}_{t\geq 0}$ es una familia creciente
de $\sigma$-álgebras: $\mathcal{F}_s \subseteq \mathcal{F}_t$ para $s \leq t$.
Intuitivamente, $\mathcal{F}_t$ representa \emph{la información disponible hasta
el tiempo $t$}. Para el movimiento browniano, la filtración natural es
$\mathcal{F}_t = \sigma(W_s : 0 \leq s \leq t)$: la $\sigma$-álgebra generada
por los valores del browniano hasta $t$.

Un proceso $X_t$ es \textbf{adaptado} a $\{\mathcal{F}_t\}$ si $X_t$ es
$\mathcal{F}_t$-medible para todo $t$: el valor de $X_t$ puede determinarse
con la información disponible hasta $t$, sin conocer el futuro. Esta condición
es la que hace que la integral de Itô tenga la propiedad de martingala: el
integrando solo puede usar información del pasado, no del futuro.

\begin{observacion}{Por qué la filtración importa en el libro}
La condición de adaptabilidad en la integral de Itô (Sección~\ref{subsec:construccion})
no es un tecnicismo vacío. Es la diferencia entre la integral de Itô (integrando
evaluado al inicio del intervalo, $\mathcal{F}_{t_i}$-medible) y la integral de
Stratonovich (integrando evaluado en el punto medio, que usa información del
futuro del intervalo). Las dos integrales difieren exactamente en el término de
corrección $\frac{1}{2}\int\partial_x f\,dt$, que es el término de Itô de la
fórmula de cambio de variables.
\end{observacion}

\subsection{Esperanza condicional como proyección $L^2$}
\label{subsec:ap_esperanza_cond}

La \textbf{esperanza condicional} $\mathbb{E}[X|\mathcal{G}]$ de una variable
$X\in L^2(\Omega)$ dada una $\sigma$-álgebra $\mathcal{G}\subseteq\mathcal{F}$
es la \emph{proyección ortogonal} de $X$ sobre el subespacio de variables
$\mathcal{G}$-medibles en $L^2(\Omega)$.

Más precisamente, $Z = \mathbb{E}[X|\mathcal{G}]$ es la única variable aleatoria
$\mathcal{G}$-medible que satisface:
\begin{equation}
  \mathbb{E}[ZY] = \mathbb{E}[XY]
  \quad\text{para toda }Y\text{ acotada y }\mathcal{G}\text{-medible}.
  \label{eq:ap_ec_proyeccion}
\end{equation}

Esta caracterización como proyección tiene consecuencias directas:

\begin{proposicion}{Propiedades de la esperanza condicional}
\label{prop:ap_ec}
\begin{enumerate}
  \item \textbf{Torre:} $\mathbb{E}[\mathbb{E}[X|\mathcal{G}]|\mathcal{H}]
        = \mathbb{E}[X|\mathcal{H}]$ para $\mathcal{H}\subseteq\mathcal{G}$.
  \item \textbf{Linealidad:} $\mathbb{E}[\alpha X + \beta Y|\mathcal{G}]
        = \alpha\mathbb{E}[X|\mathcal{G}] + \beta\mathbb{E}[Y|\mathcal{G}]$.
  \item \textbf{Extracción de lo medible:} si $Z$ es $\mathcal{G}$-medible,
        $\mathbb{E}[ZX|\mathcal{G}] = Z\,\mathbb{E}[X|\mathcal{G}]$.
  \item \textbf{Independencia:} si $X$ es independiente de $\mathcal{G}$,
        $\mathbb{E}[X|\mathcal{G}] = \mathbb{E}[X]$.
  \item \textbf{Optimalidad MSE:} $\mathbb{E}[X|\mathcal{G}]$ es el predictor
        de $X$ con la información $\mathcal{G}$ que minimiza el error cuadrático
        medio: $\mathbb{E}[(X-Z)^2]\geq\mathbb{E}[(X-\mathbb{E}[X|\mathcal{G}])^2]$
        para cualquier $Z$ $\mathcal{G}$-medible.
\end{enumerate}
\end{proposicion}

La propiedad (1) (ley de la torre) se usa implícitamente en las pruebas de
martingala del Capítulo~\ref{ch:cap10}. La propiedad (5) es la que hace de la
regresión una esperanza condicional: el score condicionado del
Capítulo~\ref{ch:cap12}, $\nabla\log p_t(\bm{x}|y)$, es el predictor MSE del
score no condicionado dado la condición $y$.

%───────────────────────────────────────────────────────────────────────────────
\section{Martingalas y convergencia}
\label{sec:ap_martingalas}
%───────────────────────────────────────────────────────────────────────────────

\textit{Usado en: Capítulo~\ref{ch:cap10} (integral de Itô como martingala,
variación cuadrática), Capítulo~\ref{ch:cap11} (proceso reverso de Anderson,
Girsanov).}

\subsection{Martingalas}
\label{subsec:ap_martingala_def}

\begin{definicion}{Martingala}
\label{def:ap_martingala}
Un proceso adaptado $\{M_t\}$ con $\mathbb{E}[|M_t|]<\infty$ para todo $t$
es una \textbf{martingala} respecto a $\{\mathcal{F}_t\}$ si:
\begin{equation}
  \mathbb{E}[M_t|\mathcal{F}_s] = M_s \quad\text{para todo }s\leq t.
  \label{eq:ap_martingala}
\end{equation}
\end{definicion}

La condición~\eqref{eq:ap_martingala} dice que el mejor predictor del valor
futuro $M_t$ dado el pasado $\mathcal{F}_s$ es el valor presente $M_s$: el
proceso no tiene drift en esperanza condicional. La imagen canónica es la de
un juego justo: las ganancias acumuladas en un juego en el que en cada paso
la ganancia esperada condicionada al pasado es cero forman una martingala.

Los tres ejemplos del Capítulo~\ref{ch:cap10} verifican este criterio:

\begin{proposicion}{Martingalas del movimiento browniano}
\label{prop:ap_mart_browniano}
Si $W_t$ es un movimiento browniano estándar, los siguientes procesos son
martingalas:
\[
  M_t^{(1)} = W_t, \quad
  M_t^{(2)} = W_t^2 - t, \quad
  M_t^{(3)} = e^{\theta W_t - \theta^2 t/2} \;\;(\theta\in\mathbb{R}).
\]
\end{proposicion}

\begin{proof}
Para $M_t^{(1)}$: $\mathbb{E}[W_t|\mathcal{F}_s] = \mathbb{E}[W_s +
(W_t-W_s)|\mathcal{F}_s] = W_s + \mathbb{E}[W_t-W_s] = W_s$, usando que
$W_t-W_s$ es independiente de $\mathcal{F}_s$ con media cero.

Para $M_t^{(2)}$: $\mathbb{E}[W_t^2-t|\mathcal{F}_s] = \mathbb{E}[(W_s+
(W_t-W_s))^2|\mathcal{F}_s] - t = W_s^2 + (t-s) - t = W_s^2 - s$.

Para $M_t^{(3)}$: $\mathbb{E}[e^{\theta W_t - \theta^2 t/2}|\mathcal{F}_s]
= e^{\theta W_s - \theta^2 s/2}\mathbb{E}[e^{\theta(W_t-W_s) - \theta^2(t-s)/2}]
= e^{\theta W_s - \theta^2 s/2}$, usando que $e^{\theta Z - \theta^2\tau/2}$
tiene esperanza 1 para $Z\sim\mathcal{N}(0,\tau)$. $\square$
\end{proof}

$M_t^{(3)}$ es la \textbf{exponencial de Wick} o \textbf{exponencial de
Doléans-Dade} en el caso de integrando determinista. En el
Capítulo~\ref{ch:cap11}, esta exponencial con integrando estocástico es la
derivada de Radon-Nikodym que aparece en el cambio de medida de Girsanov.

\subsection{Tipos de convergencia y sus relaciones}
\label{subsec:ap_convergencia}

\textit{Usado en: Capítulo~\ref{ch:cap10} (convergencia de la variación
cuadrática, construcción de la integral de Itô).}

\medskip
Sean $X, X_1, X_2, \ldots$ variables aleatorias. Hay cuatro nociones de
convergencia $X_n \to X$, de más fuerte a más débil:

\begin{enumerate}
  \item \textbf{Casi segura (c.s.):} $\mathbb{P}(\lim_{n\to\infty} X_n = X) = 1$.
        Las trayectorias convergen para casi todo $\omega$.
  \item \textbf{En $L^p$:} $\mathbb{E}[|X_n - X|^p] \to 0$.
        Para $p=2$: convergencia en media cuadrática.
  \item \textbf{En probabilidad:} $\mathbb{P}(|X_n - X| > \varepsilon) \to 0$
        para todo $\varepsilon > 0$.
  \item \textbf{En distribución:} $\mathbb{E}[f(X_n)] \to \mathbb{E}[f(X)]$
        para toda $f$ continua acotada. Solo la ley converge, no necesariamente
        las variables.
\end{enumerate}

\begin{center}
\begin{tikzpicture}[
  node distance=1.8cm,
  box/.style={draw, rounded corners, minimum width=2.2cm, minimum height=0.8cm,
    font=\small\sffamily, align=center}
]
  \node[box] (cs) {c.s.};
  \node[box, right=of cs] (lp) {$L^p$};
  \node[box, below=1.2cm of cs, xshift=1.8cm] (prob) {prob.};
  \node[box, right=of prob] (dist) {distrib.};
  \draw[->, thick] (cs) -- (prob) node[midway, left, font=\small] {};
  \draw[->, thick] (lp) -- (prob) node[midway, right, font=\small] {};
  \draw[->, thick] (prob) -- (dist);
  \node[font=\scriptsize, above=0.1cm of cs, xshift=0.9cm]
    {$\Downarrow$ ambas implican};
\end{tikzpicture}
\end{center}

Las implicaciones inversas no se sostienen en general. En el
Capítulo~\ref{ch:cap10}, la variación cuadrática converge primero en $L^2$
(cálculo de esperanza y varianza) y luego se refuerza a convergencia c.s.\ con
el lema de Borel-Cantelli.

\subsection{El segundo lema de Borel-Cantelli}
\label{subsec:ap_borel_cantelli}

El segundo lema de Borel-Cantelli es la herramienta que permite pasar de
convergencia en $L^2$ a convergencia casi segura en la demostración del
Teorema~\ref{thm:vc}. Se enuncia aquí en su forma estándar.

\begin{teorema}{Lema de Borel-Cantelli}
\label{thm:ap_bc}
Sea $\{A_n\}_{n\geq 1}$ una sucesión de eventos.
\begin{enumerate}
  \item \textbf{Primer lema:} si $\sum_{n=1}^\infty \mathbb{P}(A_n) < \infty$,
        entonces $\mathbb{P}(A_n\text{ ocurre infinitas veces}) = 0$.
  \item \textbf{Segundo lema:} si los $A_n$ son independientes y
        $\sum_{n=1}^\infty \mathbb{P}(A_n) = \infty$, entonces
        $\mathbb{P}(A_n\text{ ocurre infinitas veces}) = 1$.
\end{enumerate}
\end{teorema}

La expresión $\{A_n \text{ i.o.}\}$ (infinitely often, infinitas veces) se
escribe formalmente como $\limsup_n A_n = \bigcap_{N=1}^\infty\bigcup_{n=N}^\infty
A_n$: el conjunto de $\omega$ que pertenece a infinitos $A_n$.

En la demostración del Teorema~\ref{thm:vc} del Capítulo~\ref{ch:cap10}: se
verifica que para la subsucesión diádica $n_k = k^2$, la serie
$\sum_k\mathbb{P}(|V_T^{(2)}(\Pi_{n_k}) - T| \geq \varepsilon)$ es convergente
(de hecho, $\leq C/k^2$), y el primer lema da la convergencia c.s.\ a lo largo
de la subsucesión.

%───────────────────────────────────────────────────────────────────────────────
\section{Variables aleatorias conjuntas y condicionamiento}
\label{sec:ap_conjuntas}
%───────────────────────────────────────────────────────────────────────────────

\textit{Usado en: Capítulo~\ref{ch:cap7} (distribuciones conjuntas de pesos,
inferencia variacional), Capítulo~\ref{ch:cap12} (denoising score matching,
CFG).}

\medskip
Para variables aleatorias $(X,Z)$ con densidad conjunta $p(x,z)$:
\begin{align*}
  p(x) &= \int p(x,z)\,dz \quad \text{(marginal)}, \\
  p(z|x) &= \frac{p(x,z)}{p(x)} \quad \text{(condicional)}, \\
  p(x,z) &= p(x|z)\,p(z) = p(z|x)\,p(x) \quad \text{(factorización)}.
\end{align*}

La \textbf{ley de la esperanza total} (o ley de la torre):
\begin{equation}
  \mathbb{E}[X] = \mathbb{E}_Z[\mathbb{E}[X|Z]].
  \label{eq:ap_torre}
\end{equation}
Esta identidad se usa en la demostración del
Teorema~\ref{thm:dsm} del Capítulo~\ref{ch:cap11}: el término cruzado del score
matching se escribe como esperanza sobre $p(\bm{x}|\tilde{\bm{x}})$, que se
obtiene condicionando la esperanza conjunta sobre $p(\bm{x})\,q_\sigma
(\tilde{\bm{x}}|\bm{x})$.

\begin{observacion}{Independencia condicional y la propiedad de Markov}
Dos variables $X$ y $Y$ son \textbf{condicionalmente independientes} dado $Z$
si $p(x,y|z) = p(x|z)p(y|z)$. La \textbf{propiedad de Markov} del proceso
$\{X_t\}$ es la independencia condicional: dado $X_s$, el futuro $X_t$ ($t>s$)
es independiente del pasado $\{X_r:r<s\}$. En el Capítulo~\ref{ch:cap10},
el movimiento browniano satisface la propiedad de Markov como consecuencia
de sus incrementos independientes.
\end{observacion}

%───────────────────────────────────────────────────────────────────────────────
\begin{notasbib}
Los conceptos de este apéndice se desarrollan con mayor profundidad en los
textos estándar de probabilidad y procesos estocásticos. Para una introducción
rigurosa pero accesible a la teoría de la medida y la probabilidad: Billingsley
(1995)~\cite{billingsley1999convergence}. Para los procesos estocásticos y la
integral de Itô con un enfoque orientado a aplicaciones: Øksendal
(2003)~\cite{oksendal2003stochastic}. Para la esperanza condicional como
proyección $L^2$ y la teoría de martingalas: Karatzas y
Shreve~\cite{karatzas1991brownian}, Capítulos 1--2.

La divergencia KL y la información de Fisher en el contexto del aprendizaje
automático se tratan extensamente en Cover y Thomas
(2006)~\cite{cover2006elements}. La conexión entre la información de Fisher
y la geometría de la inferencia estadística es el tema central de Amari
(1998)~\cite{amari1998natural}.

El segundo lema de Borel-Cantelli y los tipos de convergencia se encuentran
en cualquier texto de probabilidad avanzada; la exposición aquí sigue a
Durrett (2019)~\cite{durrett2019probability}.
\end{notasbib}

\printbibliography[heading=subbibliography]